\documentclass{article}
\usepackage{/globals/anki}
\ankiprefix{infosicherheit}
\ankitopic{Informationssicherheit}
\ankishow
\begin{document}
\section{Informationssicherheit} 
Der Oberbegriff der Informationssicherheit beinhaltet unterschiedliche Aspekte, wie die Sicherstellung des Datenschutzes, der Datenintegrität und der Verfügbarkeit von Systemen und Diensten.
\ankinote{inhalt}{Inhalte}{Datenschutz, Integrität, Verfügbarkeit}

Gefahren können sowohl intern als auch extern sein.

Der Schutz der Informationssicherheit kann in drei Teile aufgeteilt werden.
\begin{description}
 \item[Prävention] Zur Prävention gehören Firewalls und Virenscanner, sinnvolle Berechtigungen, Verschlüsselungen und Awareness.
 \item[Detektion] Die Detektion passiert durch Monitoring von Abweichungen und durch Meldungen.
 \item[Reaktion] Als Reaktion sollten Systemzugriffe beschränkt, Passwörter zurückgesetzt und Behörden informiert werden.
\end{description}
\ankinote{teile}{teile}{Prävention, Detektion, Reaktion}
\ankinote{prävention}{Prävention}{Firewalls, Virenscanner, Berechtigungen, Verschlüsselung, Awareness}
\ankinote{detektion}{Detektion}{Monitoring von Abweichungen, Meldungen}
\ankinote{reaktion}{Reaktion}{Zugriffe einschränken, Passwörter zurücksetzen, Behörden informieren}
\end{document}